\silentchapter{Etterord; éin stol}{etterord; éin stol}
\begin{widebody}
{\small

Thonet nr.~14 (1859) okkuperer eitt punkt i \textbf{\textit{formrommet}}. Over femti millionar eksemplar er produserte; ho er den mest reproduserte stola i historia. Objektet tilbyr sitting, stabling og transport; ein konstellasjon av \textbf{\textit{affordansar}} som definerer klassen. \textbf{\textit{Seleksjonstrykka}} verkar samstundes; materialkostnad, produksjonshastigheit, transportvolum og sosial lesbarheit. Kvart trykk avgrensar kvar forma \textit{ikkje} kan vere. Det som står att, er nr.~14. Bøka tillèt mjuke kurver langs fiberen og vetoar sprø bøyingar; slik set materialet grensene for forma, uavhengig av den einskilde stola. Thonet navigerer etter minst fire aksar samstundes, og forma sit på ein bratt haug i \textbf{\textit{tilpassingslandskapet}} der kvart avsteg kostar. Ettertida plasserer objektet i ein kategori dei kallar historisme, men etiketten forklarer ingenting; han skildrar forma som eit ferdig fossil. Det som stivna, cellulosen, er ikkje fossil. Tinga me omgjev oss med, utgjer førebelse forlengingar av kroppane våre; protesar for muskulatur og ledd. Stola du sit i akkurat no navigerer framleis; ho responderer umerkeleg på trykket frå massen din, på rørslene til molekyla i romtemperaturen, på fotona som treff overflata. Ho er eit kompromiss mellom fysiske krefter som framleis slit i materialet. Det som skil ho frå gjenstandane bak glasmontera, er berre at oksidasjonen og friksjonen enno ikkje har fullført arbeidet sitt.

\par\addvspace{1pt}
\noindent\hfill\includegraphics[height=55mm]{thonet-14_skisse.png}\hspace{4mm}
\par\addvspace{1pt}

Tørrfurua i skogen har inga forventing om mennesket innleira i seg; ho er eit resultat av biologiske og geologiske trykk. Stola, derimot, er bunden av ei affordanse som ventar på å bli utløyst gjennom våre proporsjonar og vår tyngdekraft. Når ho står der tom, vitnar ho om eit fråvær; forma er eit negativt avtrykk, eit fysisk ekko av nokon som ikkje lenger er der. Å forme er å teikne opp grensene for eit framtidig nærvær. Kvar gjenstand er ein lovnad om menneskeleg bruk, og kvar form står att som eit fossil av ein tidlegare omhug. Alt design er til sist reiskap som ventar på eit samandrag som kanskje aldri kjem.
}
\end{widebody}
